%% ma: L11-B09-0011
%% lop: 11
%% chuong: 3
%% bai: 9
%% dang: 11.9.7
%% loai: TN4
%% muc: NB
%% dap_an: "B"
%% nguon: "de-on-tap-giua-ki-1-toan-11-nam-2026-2027-truong-thpt-ngo-quyen-da-nang.pdf (M09, Đề ôn tập giữa kì 1 Toán 11 THPT Ngô Quyền, Đà Nẵng), Mã đề 0112, Phần I Câu 8"
%% kien_thuc: "Bài 9 (các số đặc trưng đo xu thế trung tâm)"
%% trang_thai: chinh-thuc
%% ghi_chu: "Đã duyệt 29/10/2026."
\begin{ex}
Khảo sát chiều cao (đơn vị: cm) của một nhóm học sinh khối 11 thu được bảng số liệu ghép nhóm sau:
\begin{center}\small\begin{tabular}{|l|c|c|c|c|c|c|}\hline
Chiều cao & $[155;160)$ & $[160;165)$ & $[165;170)$ & $[170;175)$ & $[175;180)$ & $[180;185)$\\\hline
Số học sinh & $12$ & $25$ & $42$ & $38$ & $15$ & $8$\\\hline
\end{tabular}\end{center}
Nhóm chứa mốt của mẫu số liệu ghép nhóm trên là
\choice{$[155;160)$}{\True $[165;170)$}{$[170;175)$}{$[180;185)$}
\loigiai{Nhóm có tần số lớn nhất ($42$) là $[165;170)$.}
\end{ex}
